\documentclass[12pt,twoside,a4paper]{article}
% Ten dokument należy kompilować za pomocą XeLaTeX-a.
% Wersja dostosowana do Overleafa: zamiast lokalnego pliku arial.ttf używa fontu Liberation Sans.

\usepackage{fontspec}
\defaultfontfeatures{Ligatures=TeX}
\usepackage{graphicx}
\usepackage[vmargin=25mm,hmargin=20mm,bindingoffset=10mm]{geometry}

\newfontfamily\arial{Liberation Sans}

\linespread{1.5}
\setlength{\parindent}{0mm}

% ---------------------- Dane pracy ------------------------------

\newcommand{\discipline}{Matematyka i Analiza Danych}
\renewcommand{\title}{Wpływ funkcji pewności na jakość rekomendacji gier w modelu implicit ALS}
\renewcommand{\author}{Michał Tałałaj}
\newcommand{\album}{333092}
\newcommand{\supervisor}{dr inż. Robert Małysz}
\renewcommand{\year}{2026}

% -----------------------------------------------------------------

\begin{document}
\pagestyle{empty}

\begin{center}
\includegraphics[width=\textwidth]{img/politechnika}
\vspace{70pt}

\includegraphics[scale=1.]{img/praca_lic}

{\arial na kierunku~\discipline

\vspace{40pt}
{\arial \large \title}

\vspace{50pt}

{\arial \huge \author}

\vspace{5pt}

Numer albumu \album

\vspace{40pt}

promotor \\
{\arial \supervisor}

\vfill
WARSZAWA \year \\
}
\end{center}

\end{document}
